\section{Two saturated arithmetic directions in degree twelve}
\label{sec:twelve-normalization}

Write $E_{12}=C_\infty\cap F^{12}=E'$ for the group just obtained,
where $F^n$ denotes the dyadic word filtration. The next entire
leading layer can also be removed arithmetically.

\begin{theorem}\label{thm:twelve-normalization}
On $E_{12}$ put
\[
 \begin{split}
 b_{12,1}(g)&=[S^6TST^4]f_g,\\
 b_{12,2}(g)&=[S^5TST^2ST^2]f_g,
 \end{split}
 \qquad b_{12}=(b_{12,1},b_{12,2}).
\]
This is a continuous additive map of scalar conjugation weight
$\lambda_2^{12}$, and
\[
 b_{12}(E_{12})=b_{12}(E_{12}\cap H)=\mathbf Z_2^2.
\]
Its kernel is exactly
\[
 E_{13}=C_\infty\cap F^{13}.
\]
Consequently $E_{13}$ is a proper closed normal subgroup, with
\[
 E_{12}/E_{13}\simeq\mathbf Z_2^2,\qquad G=E_{13}H,
 \qquad
 B\simeq E_{13}/\overline{[G,E_{13}](E_{13}\cap H)}.
\]
The kernel of restriction of $X_2=H^1(N,\mathbf F_2)^G$ to
$E_{13}$ is again the odd-prime rational Kummer span plus
$\mathbf F_2\eta_3$. No nonzero linear combination of
$b_{12,1}\bmod2$ and $b_{12,2}\bmod2$ on $E_{12}$ extends to
a $G$-invariant character on $N$.
\end{theorem}

\begin{proof}
We first specify the two integral Lie directions and an exact
coefficient certificate. In lexicographically ordered Lyndon
bases, with the longest-proper-Lyndon-suffix bracketing and $S<T$,
let $L_5,L_7,L_9$ have the following integer coordinates:
\begin{align*}
L_5:\quad & (2,-4,3,4,1,-2),\\
L_7:\quad & (16,-48,80,80,-32,19,-173,-80,-32,17,-99,61,\\
          &\hspace{12mm}109,48,-65,64,48,-16),\\
L_9:\quad & (144,-576,1512,1344,-1440,-346,-6322,-2016,\\
 &360,-1598,3214,-2337,2722,4724,8823,2016,\\
 &600,2148,3259,-1118,-5303,-2951,4353,-682,\\
 &-11276,479,-1762,-8109,-6663,-1344,-89,-5496,\\
 &-1949,291,1041,7472,-4352,-2416,-3686,-1013,\\
 &2059,5340,-3340,4388,1868,3558,2434,576,\\
 &3299,-5446,-1768,6452,1368,2088,936,-144).
\end{align*}
These are the already established low-weight leading vectors,
normalized by
\[
 [S^4T]L_5=2,\qquad [S^6T]L_7=16,\qquad [S^8T]L_9=144.
\]
For any full-GT word starting in the indicated degree, its leading
term lies on the corresponding rational line. The earlier finite
leading calculations give this fact directly. It also follows from
Lemma~\ref{lem:full-gt-leading-term} and Willwacher's upper bound,
since $d_5=d_7=d_9=1$ and the known nonzero arithmetic leading
vectors supply the lower bounds. No rational Kummer normalization
is required for this statement about a word that already starts
in the specified degree.

Set
\[
 V=\{P_3,L_9\},\qquad W=\{L_5,L_7\},\qquad
 \omega_1=S^6TST^4,\quad\omega_2=S^5TST^2ST^2.
\]
Here the bracket is the same Ihara bracket used above. The exact
coefficient calculation is
\begin{equation}\label{eq:twelve-unit-minor}
 \begin{pmatrix}
 [\omega_1]V&[\omega_1]W\\
 [\omega_2]V&[\omega_2]W
 \end{pmatrix}
 =\begin{pmatrix}3385&46\\-4694&-1397\end{pmatrix},
 \qquad \det=-4512921\in\mathbf Z_2^\times.
\end{equation}
For a short reproducible certificate, the contributions of
$AB,-BA,\partial_A B,-\partial_B A$ are respectively
\[
\begin{array}{c|r|r|r|r|r}
(A,B;\omega)&AB&-BA&\partial_A B&-\partial_B A&\text{sum}\\\hline
(P_3,L_9;\omega_1)&0&0&-6720&10105&3385\\
(P_3,L_9;\omega_2)&0&-4378&12412&-12728&-4694\\
(L_5,L_7;\omega_1)&0&32&-2208&2222&46\\
(L_5,L_7;\omega_2)&0&-1152&1887&-2132&-1397
\end{array}
\]
All entries follow from the displayed Lyndon vectors by ordinary
concatenation and the finite coefficient rule
\[
 [\omega]\partial_A B=
 \sum_{B\ni w}\sum_{i:w_i=T}
       [w]B\,[\omega]\bigl(w_{<i}(TA-AT)w_{>i}\bigr).
\]
This specifies the entire integer calculation independently of
any software. In particular the two columns in
\eqref{eq:twelve-unit-minor} reduce to the identity matrix
modulo two.

We now realize odd multiples of both columns arithmetically.
Use $\sigma,h,t$ as in the previous theorem, with $t_2=3$, and
write $\tau_j=\tau_j(\sigma)$ and $a=m_3(\sigma)$; all of
$a,\tau_5,\tau_7,\tau_9$ are odd. Put
\[
 u_5=\Delta_3(\sigma),\qquad
 u_7=\Delta_5(u_5),\qquad
 \widetilde u_9=\Delta_7(u_7),\qquad
 \Delta_j(z)=hzh^{-1}z^{-t^j}.
\]
Every exponent $t^j$ is the full profinite multiplier power.
The operator $\Delta_j$ kills the homogeneous leading term of
degree $j$. There is no full-GT leading term in degrees four or
six: Lemma~\ref{lem:full-gt-leading-term} and the bound $d_4=d_6=0$
give this immediately. Thus $u_5\in F^5$, $u_7\in F^7$, and
$\widetilde u_9\in F^8$.

The degree-eight leading space is the single line detected by
$m_8$, and $w_8$ has an odd $m_8$ value. Choose
$r\in\widehat{\mathbf Z}$ with
$r_2=m_8(\widetilde u_9)/m_8(w_8)$ and put
\[
 u_9=\widetilde u_9 w_8^{-r}.
\]
Then $u_9\in F^9\cap H$. Since $B_{w_8}=1$, this correction
does not change any beta moment.

The additive weight laws for the moments and the triangular
depth-one formulas give the exact leading terms
\begin{equation}\label{eq:twelve-primitive-inputs}
 \begin{split}
 \operatorname{in}_5(u_5)&=9\tau_5 L_5,\\
 \operatorname{in}_7(u_7)&=729\tau_7 L_7,\\
 \operatorname{in}_9(u_9)&=2302911\tau_9 L_9.
 \end{split}
\end{equation}
Indeed, after the preceding moments have been killed, the raw
depth-one coordinate in degree $2k+1$ is
$2\tau_{2k+1}/(2k)!$. Therefore the three raw coordinates are
\[
 \begin{split}
 m_5(u_5)&=(3^5-3^3)\tau_5/12=18\tau_5,\\
 m_7(u_7)&=(3^7-3^5)(3^7-3^3)\tau_7/360=11664\tau_7,\\
 m_9(u_9)&=(3^9-3^7)(3^9-3^5)(3^9-3^3)\tau_9/20160\\
         &=331619184\tau_9.
 \end{split}
\]
Dividing these by the displayed depth-one coefficients $2,16,144$
gives \eqref{eq:twelve-primitive-inputs}. Only odd scalars
multiply the primitive integer vectors. This construction does
not divide an arithmetic group element by two.

Define
\[
 z_{3,9}=[\sigma,u_9],\qquad z_{5,7}=[u_5,u_7].
\]
These lie in $H\cap C_\infty\cap F^{12}=H\cap E_{12}$:
their leading degrees are at least $3+9$ and $5+7$, and every
rational Kummer coordinate and the beta multiplier vanish on
commutators in $M$. Their leading terms are
\[
 \operatorname{in}_{12}(z_{3,9})=2302911a\tau_9 V,
 \qquad
 \operatorname{in}_{12}(z_{5,7})=6561\tau_5\tau_7 W.
\]
Thus both columns of \eqref{eq:twelve-unit-minor} are realized
up to odd arithmetic factors.

Willwacher's Theorem~1.1 and the base-change lemma give
$\dim\mathfrak{grt}_1(\mathbf Q_2)_{12}\le d_{12}=2$.
To count $d_{12}$ directly, the only nonzero bracket types are
$(3,9)$ and $(5,7)$: four weight-three letters span zero in a
free Lie algebra on that single letter, and all other generator
lengths or weights are excluded. The two arithmetic leading
vectors just constructed are independent by the odd determinant,
so they span the full rational leading space in degree twelve.

Every leading polynomial $q$ of an element of $E_{12}$ is
therefore $cV+dW$ with $c,d\in\mathbf Q_2$. Its two selected
Magnus coefficients are integral. The inverse of the matrix in
\eqref{eq:twelve-unit-minor} is integral over $\mathbf Z_2$,
so $c,d\in\mathbf Z_2$. Conversely, products of profinite powers
of $z_{3,9},z_{5,7}$ realize every such pair, since their scalar
factors are odd. Hence the entire integral leading lattice, and
its arithmetic image, is $\mathbf Z_2V+\mathbf Z_2W$.

Leading coefficients on $F^{12}$ add under the GT product and
have scalar conjugation weight $\lambda_2^{12}$. The unit minor
shows that vanishing of $b_{12}$ is equivalent to vanishing of
the whole degree-twelve term. This proves the asserted kernel
and image statements. Arithmetic correction by the two words
gives $E_{12}=E_{13}(E_{12}\cap H)$, whence $G=E_{13}H$.
The same closed normal subgroup argument used in the strict
degree-eleven reduction identifies the relative quotient $B$.

Finally the two correcting words are killed by every original
$N$-invariant character. To see this with all coefficient primes
preserved, put
\[
 c=(t-t^7)(t-t^5)(t-t^3)\in\widehat{\mathbf Z}.
\]
All rational Kummer coordinates satisfy
$\rho_q(u_9)=c\rho_q(\sigma)$; the degree-eight correction has
zero rational Kummer coordinates. Thus $u_9\sigma^{-c}\in D$
and
\[
 z_{3,9}=[\sigma,u_9\sigma^{-c}].
\]
Likewise $\rho_q(u_7)=(t-t^5)\rho_q(u_5)$ for every $q$, so
$u_7u_5^{-(t-t^5)}\in D$ and
\[
 z_{5,7}=[u_5,u_7u_5^{-(t-t^5)}].
\]
Both are commutators $[g,n]$ with $g\in G$ and $n\in D\subset N$.
Every character in $X_2$ therefore vanishes on both. Correcting
by their profinite powers shows that its restriction to $E_{13}$
vanishes if and only if its restriction to $E_{12}$ vanishes.
This proves the unchanged kernel. The same two evaluations,
whose matrix is invertible modulo two, exclude every nonzero
linear combination of the two quotient characters as an
invariant restriction from $N$.
\end{proof}

\begin{remark}
Both new quotient directions are already arithmetic. Their
removal preserves the relative obstruction; it does not exhibit
a nonzero element of $\operatorname{Hom}(B,\mathbf F_2)$.
The proof uses a rational dimension upper bound together with
an explicit odd integral minor. Neither ingredient substitutes
for the other.
\end{remark}
